% =========================================================================
% MSc Dissertation — opening chapters (draft 1, 2026-08-07)
% Grounded in the codebase as handed over (branch `handover`); every claim
% about the system below is traceable to a module named in the text.
% Placeholders are marked \todo{...} — search for "todo" before submitting.
% =========================================================================
\documentclass[11pt,a4paper]{report}

\usepackage[margin=2.5cm]{geometry}
\usepackage[hidelinks]{hyperref}
\usepackage{booktabs}
\usepackage{xcolor}
\usepackage{parskip}

\newcommand{\todo}[1]{\textcolor{red}{[TODO: #1]}}

\title{\todo{Title} \\[1ex]
  \large A source-verified entity-extraction pipeline for consultancy company research}
\author{George Stamatopoulos \\
  \todo{Department}, Imperial College London \\
  In collaboration with Sagentia Innovation \\
  \texttt{gs2423@ic.ac.uk}}
\date{\todo{Submission date}}

\begin{document}

\maketitle

\begin{abstract}
Technology and product-development consultancies routinely begin engagements
with a company-research phase: given a list of candidate companies and a set
of client questions, analysts visit each company's website and fill in an
answer matrix by hand. This dissertation presents a pipeline, built during an
industrial placement with Sagentia Innovation, that automates this phase
while remaining auditable by construction: every extracted answer carries a
verbatim quote from its source page, and an independent verification layer
checks that the quote genuinely appears there --- the language model is never
taken at its word. The system was developed against the constraints of a real
corporate deployment environment (locked-down machines, no model downloads,
strict crawling politeness) and evaluated on a live advisory engagement:
59 contract-manufacturing organisations, 17 client questions, with ground
truth produced independently by three human analysts. On this benchmark the
pipeline reaches an answer-level F1 of 0.606 (precision 0.596, recall 0.617),
and a downstream experiment shows that a criteria-driven shortlist computed
from the pipeline's output agrees with the analyst-derived shortlist on four
of the top five companies. The dissertation analyses where and why the system
succeeds and fails, and what those failure classes imply for deploying
LLM-based extraction in professional advisory work.
\end{abstract}

\tableofcontents

% =========================================================================
\chapter{Introduction}
% =========================================================================

\section{Context and motivation}

Sagentia Innovation is a science and technology advisory business. A
recurring early step in its engagements is \emph{company screening}: a client
supplies a list of candidate companies --- suppliers, competitors, acquisition
targets, contract manufacturers --- together with the questions that matter to
them (\emph{Where do they manufacture? What is their revenue? Do they have
medical-device experience?}), and analysts answer those questions from the
public evidence, principally each company's own website. The output is an
answer matrix: one row per company, one column per question, every cell
ideally backed by something the analyst actually read.

Done by hand, this work is slow, repetitive, and unevenly thorough: an
analyst facing 100 companies and 15 questions makes 1,500 judgement calls,
each requiring navigation of a different website. It is also exactly the kind
of reading-and-transcribing task that large language models appear well
suited to --- and exactly the kind of task where their well-documented failure
mode, confident fabrication, is commercially unacceptable. A consultancy
cannot put a hallucinated revenue figure in front of a client.

The engineering question of this dissertation is therefore not \emph{can an
LLM answer questions about companies?} --- it evidently can --- but \emph{can a
system be built around an LLM such that its output is trustworthy enough,
and demonstrably enough so, to enter a professional deliverable?}

\section{The problem}

Concretely, the system must accept a consultant-shaped input --- an Excel
workbook listing entities, seed URLs, and free-text questions --- and produce
a consultant-shaped output: the filled answer matrix, plus the evidence
trail that lets a sceptical reviewer check any cell. Between input and
output lie several distinct hard problems:

\begin{enumerate}
  \item \textbf{Acquisition.} Company websites vary from static brochure
    sites to JavaScript-rendered single-page applications behind bot
    detection. Relevant evidence may sit several links deep, and crawl
    budget spent on privacy policies is budget not spent on factory-location
    pages.
  \item \textbf{Extraction.} Answers must be pulled from noisy marketing
    prose into structured values --- with the distinction preserved between
    \emph{the site says no}, \emph{the site does not say}, and \emph{we
    could not read the site}.
  \item \textbf{Grounding.} Each answer must be tied to verbatim source
    text, and that tie must be checked mechanically, not asserted by the
    same model that produced the answer.
  \item \textbf{Aggregation and presentation.} Evidence from many pages must
    collapse to one defensible answer per cell, without hiding genuine
    conflicts between sources.
  \item \textbf{Evaluation.} ``Does it work?'' must be answerable with
    numbers, against ground truth produced by the people whose work the
    system is meant to support --- including an analysis of whether its
    errors would change the decisions a consultant makes downstream.
\end{enumerate}

A further, less conventional set of constraints shaped the work throughout:
the system had to run on corporate IT infrastructure where model downloads
are blocked, new package installations are negotiated case by case, and
crawling misbehaviour has previously led to company IP addresses being
blocked. Several architectural decisions that would look eccentric in a
research prototype --- a dependency-free web interface, embeddings served only
from an internal server with a lexical fallback, a static-first fetcher that
escalates to browser rendering only on demonstrated need --- are direct
consequences of designing for deployment rather than demonstration.
Where elegance and deployability conflicted, deployability won.

\section{Approach and contributions}

The system is a layered pipeline (Chapter~\ref{ch:system}) in which the LLM
is deliberately confined: it reads pages and proposes
\{\emph{value}, \emph{verbatim quote}\} pairs, and everything else ---
crawling, routing, quote verification, deduplication, conflict flagging,
ranking, shortlisting --- is deterministic code that can be tested, replayed,
and audited. The dissertation's contributions are:

\begin{enumerate}
  \item \textbf{A production-shaped, source-verified extraction pipeline}
    (Chapter~\ref{ch:system}): five core layers (Acquire, Filter, Extract,
    Verify, Aggregate) with an optional cited-summary layer whose every
    sentence must reference a verified claim, enforced by a mechanical gate
    with a verbatim fallback.
  \item \textbf{An evaluation framework and ground-truth methodology}:
    a generic evaluator aligning pipeline output to analyst-filled ground
    truth via a cross-encoder matcher with typed deterministic rules for
    numerics and binaries; three synthetic benchmark tasks; and a
    59-company, 17-question ground truth built independently by three
    analysts on a live engagement.
  \item \textbf{An empirical analysis of failure classes}: which errors come
    from acquisition (starved crawls, bot walls, JavaScript-hidden
    navigation), which from extraction (evidence-strength judgement,
    entity attribution), and which from a genuine norm boundary between
    what a website supports and what an analyst is willing to assert ---
    together with measured evidence on which of these respond to
    engineering effort and which do not.
  \item \textbf{A decision-level evaluation}: a criteria-driven shortlist
    layer, with criteria expressed as data rather than code, used to
    measure whether pipeline-vs-analyst disagreements at the cell level
    actually change the companies a consultant would shortlist.
\end{enumerate}

\section{Report structure}

Chapter~\ref{ch:system} describes the system as built, layer by layer, and
the design principles that connect them. \todo{Subsequent chapters: the
evaluation framework and ground-truth methodology; the case-study results
and failure-class analysis; the shortlist and decision-level analysis;
discussion, limitations and future work. Fix chapter numbering when those
chapters exist.}

% =========================================================================
\chapter{System overview}
\label{ch:system}
% =========================================================================

\section{Design principles}

Four principles recur in every layer and explain most of the architecture:

\paragraph{The LLM is never trusted on its own word.} Every extracted answer
must arrive with a verbatim quote, and a deterministic verification layer
independently confirms the quote appears in the cached source page. The same
principle extends to the optional AI summary: each sentence must cite a
claim ID from the verified set, and a mechanical gate replaces any summary
that fails to comply with the verbatim claims themselves. Trust is
established by construction, not by model choice.

\paragraph{Determinism wherever possible.} Only two stages call a generative
model (Extract, and the optional Summarize); both run at temperature zero
with a fixed seed. Everything else --- link scoring, routing, quote checking,
deduplication, conflict detection, claim-ID assignment, shortlist scoring ---
is deterministic, so runs are reproducible, failures are replayable, and
A/B comparisons measure the change under test rather than sampling noise.

\paragraph{Honest abstention.} The pipeline distinguishes \emph{the site
states there is none}, \emph{the site does not say}, and \emph{no usable
answer found}, and these sentinels are never allowed to compete with real
answers during aggregation. An empty cell is a finding, not a blank: the
end-of-run summary names the sites that yielded nothing and why (e.g.\
bot-wall refusals), because for screening questions such as ``is this
company still independent?'', absence of evidence is itself evidence a
consultant needs to see.

\paragraph{Deployability over elegance.} No component may assume
capabilities the target environment lacks. Embeddings come from an internal
Ollama server (\texttt{nomic-embed-text}) with a BM25 lexical fallback when
it is unreachable; the user interface is a local web page served by the
Python standard library rather than a JavaScript build stack; crawling
enforces robots.txt, per-domain delays, and an honest user agent by
construction.

\section{Pipeline architecture}

The pipeline (orchestrated in \texttt{pipeline.py}, stages under
\texttt{src/}) reads an input workbook of entities, seed URLs, questions,
and optional per-run configuration, and executes the following stages per
entity, concurrently across entities:

\paragraph{Acquire (\texttt{src/acquire/}).} A static-first hybrid fetcher
(\texttt{fetcher.py}): a plain HTTP request with content extraction
(Trafilatura), judged by a quality gate; only on gate failure does it
escalate to a headless Chromium render (Playwright), keeping whichever text
is richer. A guided same-domain crawler (\texttt{crawler.py}) scores
discovered links against the questions (embedding similarity, BM25
fallback) and follows the most promising within a page budget; scope and
render-assisted discovery are configurable per run. All pages are cached by
URL hash, making re-runs cheap and experiments replayable.

\paragraph{Filter (\texttt{src/filter.py}).} Routes each fetched page to the
questions it could plausibly answer, by embedding-similarity threshold or
keyword match. The layer is fail-safe by contract: it never discards a
page --- at worst it routes all questions --- so it can reduce cost but never
recall.

\paragraph{Extract (\texttt{src/extract.py}).} The single load-bearing LLM
stage. Production uses Azure OpenAI GPT-4.1-mini (the backend is pluggable),
prompted for strict JSON \{\emph{value}, \emph{verbatim quote}\} pairs per
entity and question, over overlapping page chunks, at temperature zero with
a fixed seed and per-chunk caching. The prompt forbids extracting absence
markers as if they were answers; explicit on-page non-disclosure statements
remain extractable.

\paragraph{Verify (\texttt{src/verify.py}).} Deterministic quote grounding:
each quote is checked against the cached page text --- exact substring, then
fuzzy match, then soft anchors for long quotes --- plus a diagnostic semantic
score. Unverified evidence is retained but flagged and demoted, never
silently dropped: the reviewer sees what failed verification and why.

\paragraph{Aggregate (\texttt{src/aggregate.py}).} Collapses per-page
evidence to per-cell answers: fuzzy deduplication of near-identical values,
conflict flagging on single-answer questions (conflicts are displayed
side by side, never merged away), evidence ranked by verification quality.
For single-answer cells the rendered matrix leads with the best-supported
candidate and points to the full evidence sheet for the rest.

\paragraph{Group and Summarize (\texttt{src/group.py},
\texttt{src/summarize.py}).} Verified claims are clustered into themes and
assigned stable citable IDs (\texttt{[C\#\#\#\#]}), forming the audit trail.
An optional summary layer renders each cell for reading --- short factual
cells deterministically, multi-value cells through a citation-pooling merge,
prose cells through an LLM synthesis in which every sentence must cite claim
IDs; a mechanical gate falls back to verbatim claims whenever a sentence
fails to comply. Both layers are strictly additive and fail-soft: their
failure cannot fail a run.

\paragraph{Output (\texttt{src/io\_excel.py}).} A single Excel workbook:
the answer Matrix and Summary for consumption; Provenance, Digest and
Grouped Themes for audit; per-stage diagnostic logs for engineering. Every
number a consultant might quote is at most one sheet away from the verbatim
text it came from.

\paragraph{Downstream: shortlist (\texttt{src/shortlist.py}).} A post-hoc,
read-only layer that turns the finished matrix into a ranked shortlist under
client criteria expressed entirely as data (a criteria workbook: gates,
thresholds, weights, keyword lists, missing-data policies). Its parsers are
deterministic and refuse to guess --- ambiguous values fail with a stated
reason and every assumption becomes a visible flag --- which makes the layer
usable in reverse, as an instrument: running identical criteria over the
pipeline's matrix and over analyst ground truth measures, at the level of
decisions rather than cells, how much the pipeline's errors actually matter.
This experiment is reported in Chapter~\todo{ref}.

\section{Evaluation harness}
\label{sec:eval-harness}

The evaluation framework (\texttt{src/eval/}) is code of the same standing
as the pipeline: a generic evaluator aligns pipeline cells to ground-truth
cells and classifies each pairing (match, miss, false assertion, null
agreement), using a cross-encoder semantic matcher wrapped in typed
deterministic rules --- numerics and binaries are compared by parsed value
and polarity, because semantic similarity models mis-handle exactly those
(\emph{Yes} vs \emph{No} score as similar; \emph{6{,}500+} vs \emph{6500+}
as dissimilar). Ground truth is authored by analysts in the same template
the pipeline fills, under a written convention (website-only evidence;
\emph{No} requires positive evidence; ties go to \emph{Not disclosed}) ---
the methodology, its inter-annotator checks, and the resulting benchmark
are the subject of Chapter~\todo{ref}.

\todo{End of drafted material. Next chapters: Evaluation framework \&
ground truth methodology; Case study \& failure-class analysis; the
decision-level shortlist experiment; Discussion \& limitations.}

\end{document}
